% SPDX-FileCopyrightText: 2026 Will Cook
% SPDX-License-Identifier: CC-BY-4.0
% System-architecture paper; build with tectonic.
\documentclass[10pt]{article}
\usepackage{paper-house-style}
\usepackage{booktabs,array,float,placeins,multicol}
\usepackage{tikz}
\usetikzlibrary{arrows.meta,positioning,calc,fit,shapes.geometric}
\usepackage[colorlinks=true,linkcolor=linkinternal,urlcolor=linkexternal,citecolor=linkinternal]{hyperref}
\hypersetup{
  bookmarksnumbered=true,bookmarksopen=true,bookmarksopenlevel=1,
  pdftitle={Problem-Sized Lean Worlds},
  pdfauthor={Will Cook},
  pdfsubject={Architecture and trust boundaries of an AI-assisted Lean research system},
  pdflang={en-GB},
  pdfkeywords={formal mathematics, Lean 4, AI agents, research architecture, claim verification, reproducibility},
}
\newcommand{\repobase}{https://github.com/wcook04/plectis-erdos}
% BEGIN generated_semantic_coverage_macros
\newcommand{\SemanticAuthoredTheoremLike}{144,487}
\newcommand{\SemanticAuthoredInterpreted}{139,799}
\newcommand{\SemanticDirectEvidence}{3,311}
\newcommand{\SemanticContextual}{136,488}
\newcommand{\SemanticStructuralOnly}{4,688}
\newcommand{\SemanticAuthoredPercent}{96.8}
% END generated_semantic_coverage_macros
\newcommand{\repolink}[2]{\href{\repobase/blob/ca0e13f8acf5ccf48506e4bdb870953d3a0856fa/#1}{\texttt{#2}}}
\emergencystretch=1.7em

\runninghead{Problem-sized Lean worlds}
\title{Problem-Sized Lean Worlds}
\subtitle{How AI-assisted research is recorded, checked, and published}
\author{Will Cook}
\date{September 2026}

\begin{document}
\maketitle
\paperprovenance
\fontsize{9.2pt}{11pt}\selectfont

\begin{abstract}
Plectis is a research environment for studying open mathematical problems with
AI agents. It records the statements already proved, the approaches that
failed, and the questions still worth pursuing, together with the source
material needed to examine them. The aim is to let a new research attempt
begin where the previous work stopped.

This paper describes the implementation. Agents retrieve a small relevant
part of the research record, save their proof attempts and check results,
and update the formal source and papers when a result changes. A public Git
repository makes the published work available for others to read, reproduce,
and extend. We explain the design through a failed extension of a
multiplication identity and a finite certificate whose scope must be preserved
in publication. The contribution is the integration of these research and
publication operations in a working prototype. It covers eight Erd\H{o}s
problems, all still open. The empirical material reported here is of two
kinds. An internal agent rehearsal is a run of the author's own agents inside
this repository. A deterministic mutation test seeds a known false edit and
records whether a configured check rejects it. Neither kind measures how an
outside reader understands the corpus. The implementation and author-operated
checks are available; whether this organisation improves research productivity
has yet to be tested with outside contributors.
\end{abstract}

\section{The problem: many kinds of evidence}

A researcher returning to a difficult problem needs to know what has changed.
Which statements have been proved? Why did the last approach stop? Has someone
already tried the calculation now being proposed? A finished paper answers
some of these questions, but much of the work can disappear between attempts.
That loss is particularly easy when an AI agent's context ends and another
agent starts from a summary.

We built Plectis to keep this work available. Each problem has a versioned
record containing its literature, formal statements, computations, unsuccessful
approaches, and papers. An agent can query that record, investigate a smaller
question, and save a result that another researcher can inspect. The public
repository contains the published mathematics and the tools for continuing it;
the private workbench coordinates the agents that produced the initial work.
The \papersectionlink{open-source-mathematics-strategy.pdf}{strategy-protocol}%
{companion strategy paper} explains how outside contributions return through
GitHub with evidence and credit.

\paragraph{What this paper contributes.}
We describe how the implementation connects research records to formal checks
and public explanations. Lean, the proof assistant used here, checks a proof
of a precisely stated proposition. A separate record links selected
propositions to the wording approved for the paper. Release checks then test
those recorded relationships when files change. A change to a proof's
hypotheses can therefore trigger a review of the explanation that depends on
it, rather than leave the paper describing an older statement.

Linking mathematical statements to named Lean declarations is established
prior art. \texttt{leanblueprint} already records the declarations that
realise a statement, the statement and proof dependency edges an author
declares, and a formalisation-readiness state for each node, and its
declaration check confirms that every named declaration exists in the project
or one of its dependencies \cite{leanblueprint}. This repository claims none
of that layer as new. The object maintained here is the relationship among the
public mathematical claims, the several evidence classes that support them
(Lean-checked, ordinary proof, and computational), the failed routes whose
obstructions have themselves been proved, the generated public projections a
reader navigates, and the research task that continues after the document is
published.

Several parts have close precedents. Agent Hunt coordinates formalisation
work, LeanMarathon maintains a shared proof blueprint across long agent runs,
and Lean Atlas narrows the declarations that require semantic inspection
\cite{agenthunt,leanmarathon,leanatlas}. Our contribution is architectural
composition around a continuing problem record. Tao's account of proof
generation, verification, exposition, and digestion motivates including the
explanatory work in that record \cite{taoai}. The comparison in
Section~\ref{sec:related} identifies the overlaps.

\subsection{What the system records}

We use \emph{problem world} for the collection associated with one open
problem: its statement, literature, partial results, experiments, and failed
approaches. A query selects a bounded neighbourhood inside a problem-sized
world. For example, an attempt to weaken a theorem's hypothesis needs that
theorem, the declarations its proof uses, and known counterexamples to the
weaker claim. It rarely needs the entire library.

The stages require different checks. Concurrent agents can overwrite each
other's files. A proof can establish a proposition that was stated incorrectly.
A paper can exaggerate a correctly proved result. Table~\ref{tab:nonfungible}
shows how the implementation handles these different failures.

\begin{table}[!ht]
\centering\scriptsize
\begin{tabular}{@{}>{\raggedright\arraybackslash}p{0.19\linewidth}
                    >{\raggedright\arraybackslash}p{0.30\linewidth}
                    >{\raggedright\arraybackslash}p{0.39\linewidth}@{}}
\toprule
\papertablehead
Operation & Record & Check \\
\midrule
Select a question & Relevant statements and source references & Inspect the selected context and its omissions. \\
Edit shared files & Time-limited claim on named paths & Reject overlapping work claims. \\
Run a proof check & Source revision and build result & Use the pinned Lean environment and available build slot. \\
Describe a result & Statement, evidence, and hypotheses & Reconcile the formal and informal claims. \\
Publish a revision & Reviewed wording and source links & Check the registered publication relationships. \\
Request external review & Selected results and review outcome & Record the review separately from local acceptance. \\

\bottomrule
\end{tabular}
\caption{Checks at different stages of the research process.}
\label{tab:nonfungible}
\end{table}

The system also records repairs to its own tools. A missing source link may
lead to a better query; a repeated build conflict may lead to a scheduling
change. These are changes to the research process. Mathematical records change
when a proof, counterexample, computation, or reviewed statement supplies the
corresponding evidence.

\section{The whole lifecycle in one picture}\label{sec:lifecycle}
\papersectiontarget{systems-lifecycle}

Figure~\ref{fig:lifecycle} follows an attempt from question selection to
publication. A failed proof returns a specific unfinished obligation to the
researcher. A reviewer may instead find a misstatement or an overlooked
source. The next attempt starts with that correction recorded.

\begin{figure}[!t]
\centering
\begin{tikzpicture}[
  box/.style={paper stage, text width=2.05cm},
  private/.style={box, fill=MidnightBlue!5, draw=MidnightBlue!55},
  formal/.style={box, fill=TealBlue!6, draw=TealBlue!60!black},
  public/.style={box, fill=ForestGreen!6, draw=ForestGreen!55!black},
  assure/.style={box, fill=BurntOrange!7, draw=BurntOrange!65!black},
  arr/.style={paper arrow},
  back/.style={paper feedback},
  lab/.style={font=\scriptsize\sffamily\bfseries, text=gray!70!black}
]
\node[private] (intent) at (0,0) {{\scriptsize\sffamily\bfseries 1 \; SELECT}\\[3pt]
  \textbf{Intent and route}\\[2pt]\scriptsize question, scope, claim};
\node[private, right=3.5mm of intent] (reason) {{\scriptsize\sffamily\bfseries 2 \; TEST}\\[3pt]
  \textbf{Reason and test}\\[2pt]\scriptsize candidates and failures};
\node[formal, right=3.5mm of reason] (lean) {{\scriptsize\sffamily\bfseries 3 \; CHECK}\\[3pt]
  \textbf{Formalise}\\[2pt]\scriptsize proof and kernel verdict};
\node[public, right=3.5mm of lean] (claim) {{\scriptsize\sffamily\bfseries 4 \; EXPLAIN}\\[3pt]
  \textbf{Public claim}\\[2pt]\scriptsize result, evidence, limit};
\node[assure, right=3.5mm of claim] (review) {{\scriptsize\sffamily\bfseries 5 \; RELEASE}\\[3pt]
  \textbf{Assure}\\[2pt]\scriptsize replay and boundary checks};
\draw[arr] (intent) -- (reason);
\draw[arr] (reason) -- (lean);
\draw[arr] (lean) -- (claim);
\draw[arr] (claim) -- (review);
\draw[back] (review.south) to[out=-115,in=-65,looseness=1.15]
  node[lab, below=4pt] {record what failed or needs revision} (reason.south);
\node[font=\scriptsize\sffamily\bfseries, text=MidnightBlue!75!black,
  above=5mm of $(intent.north)!0.5!(reason.north)$] {PRIVATE WORKBENCH};
\node[font=\scriptsize\sffamily\bfseries, text=ForestGreen!62!black,
  above=5mm of $(lean.north)!0.5!(review.north)$] {SELF-CONTAINED PUBLIC ARTEFACT};
\draw[paper boundary] ($(reason.east)!0.5!(lean.west)+(0,-9mm)$) --
  ($(reason.east)!0.5!(lean.west)+(0,12mm)$);
\end{tikzpicture}
\caption{The end-to-end architecture. The dashed vertical line marks the
release boundary. Everything to its right remains usable without the private
workbench.}
\label{fig:lifecycle}
\end{figure}

The agent selects a question and gathers the statements and evidence needed
for it. It then runs an experiment, edits a proof, or writes an explanation.
The result and its checks are saved with the source version, so another attempt
can use them. A finding may also suggest a reusable lemma, a better query, or
a repair to the tools.

The public proof-cockpit command summarises the checkout, Lean toolchain,
open propositions, and saved research sessions. Its check mode runs the
claim-registry, cold-clone, and projection-freshness checks. These establish
whether the repository records agree with their sources; checking a proof
requires a separate Lean run. The commands are listed in
Appendix~\ref{app:repro}.

\subsection{Instructions for recurring jobs}
\papersectiontarget{systems-job-lifecycle}

The public clone provides a skill file for each recurring job. It tells an
agent where to start, what it may change, what evidence to save, and when to
stop or hand the work on. Figure~\ref{fig:job-lifecycle} shows the ordinary
research path.

\begin{figure}[!t]
\centering
\begin{tikzpicture}[
  job/.style={paper stage, text width=2.35cm, minimum height=1.25cm},
  arr/.style={paper arrow}
]
\node[job, fill=MidnightBlue!5] (invoke) at (0,0)
  {{\scriptsize\sffamily\bfseries 1}\\[2pt]\textbf{Invoke}\\[2pt]\scriptsize bounded route};
\node[job, fill=MidnightBlue!5, right=5mm of invoke] (work)
  {{\scriptsize\sffamily\bfseries 2}\\[2pt]\textbf{Work}\\[2pt]\scriptsize continue or suspend};
\node[job, fill=TealBlue!5, right=5mm of work] (validate)
  {{\scriptsize\sffamily\bfseries 3}\\[2pt]\textbf{Validate}\\[2pt]\scriptsize computation or Lean};
\node[job, fill=TealBlue!5, right=5mm of validate] (propagate)
  {{\scriptsize\sffamily\bfseries 4}\\[2pt]\textbf{Propagate}\\[2pt]\scriptsize classify consequences};
\node[job, fill=ForestGreen!6, below=8mm of propagate] (package)
  {{\scriptsize\sffamily\bfseries 5}\\[2pt]\textbf{Package}\\[2pt]\scriptsize evidence and provenance};
\node[job, fill=ForestGreen!6, left=5mm of package] (propose)
  {{\scriptsize\sffamily\bfseries 6}\\[2pt]\textbf{Propose}\\[2pt]\scriptsize pull request or issue};
\node[job, fill=BurntOrange!7, left=5mm of propose] (adopt)
  {{\scriptsize\sffamily\bfseries 7}\\[2pt]\textbf{Review and adopt}\\[2pt]\scriptsize reconcile, accept, credit};
\draw[arr] (invoke) -- (work);
\draw[arr] (work) -- (validate);
\draw[arr] (validate) -- (propagate);
\draw[arr] (propagate) -- (package);
\draw[arr] (package) -- (propose);
\draw[arr] (propose) -- (adopt);
\end{tikzpicture}
\caption{A contribution from the initial task to its inclusion in the repository.}
\label{fig:job-lifecycle}
\end{figure}

The public skills are instructions an agent can execute in a clone. They cover
orientation, research, proof checks, updating affected papers and records, and
preparing a pull request. The reproducibility appendix and the agent entry
link to the individual commands.

\section{The private workbench}\label{sec:private}

\subsection{Disk is shared memory}

Durable state lives in files. The workbench stores the task, source changes,
check results, and unfinished work so another agent can resume without relying
on the previous conversation. A \emph{receipt} is a saved record of an
operation and its outcome; a \emph{projection} is a generated view of the
underlying records, such as an index or status page.

A router selects a short set of relevant records and instructions for the
current task. Each summary points to its source, which the agent can open when
it needs more detail. The mathematical query mechanism is described in
Section~\ref{sec:mathloop}.

\subsection{Type A and Type B}

Type A is an agent working in a harness such as Claude Code or Codex. It can
read the checkout, run tools, edit files, and test the result. Type B is a
one-shot LLM exchange outside that harness, such as Astra Pro in the web app.
The operator supplies a research packet and brings the response back into the
checkout, where the harness agent checks and integrates it. The labels describe
substrate access, not model quality.

The public \texttt{source\_packet.py} tool exports selected committed files
into one readable attachment. It records the source commit and file hashes,
preserves the complete selected text, and offers an offline checksum check.
An operator can therefore send the same statement and surrounding evidence
to a web-app model without asking it to navigate GitHub. The operator chooses
the sources; the tool does not infer their dependency closure. Binary files
are preserved as labelled base64 and require decoding before inspection.

The private packet builder also prepares for two ways of reading. A browser-only
recipient receives links to the public frontier and its source statements. An
offline dossier carries the selected proofs and their dependencies. When it
includes a runnable checker, the builder also packages its declared local
runtime dependencies and rejects missing ones. This lets a recipient with code
execution test the supplied calculation without reconstructing the author's
checkout. The returned answer still goes through local checking and integration.

\subsection{Concurrency without shared-state confusion}

Several agents may work at once, but concurrency is admitted only where the
write scopes can be separated. A work item identifies the objective and
expected evidence. Before mutation, an agent claims exact paths for a bounded
lease. Independent claims may proceed concurrently; overlapping claims must be
coordinated or deferred. Fan-out is followed by a fan-in barrier where the
results are compared, validated, and integrated.

Coordination state is written to append-only ledgers and immutable receipts;
generated status pages project those records. Bounded job counts and focused
Lean builds keep concurrent work within the machine's capacity.

\subsection{Routing and scheduling}

The router exposes a typed option surface: searchable entries for tools,
instructions, source files, and ongoing work. An agent first receives a short
card, then follows its links to the working context or source. The selected
entries remain visible, so a missing premise or poor retrieval choice can be
investigated. The initial response gives the next action and a route to the
full record. Detailed audits stay available without being repeated at every
entry; task-specific constraints remain in the selected context.

For a mathematical query, the compiler assembles the target statement and
premises for possible arguments. Feedback revises the affected argument while
retaining earlier counterevidence and diagnostics. The packet records source
fingerprints; a freshness check compares them with the current inputs and
requests recompilation if they differ.

The work ledger separates past work from present permission to edit. An agent
claims exact paths for a bounded lease. A directory claim conflicts with claims
on its children, and a lease expires even if its holder crashes. Completion,
reopening, and replacement are recorded as separate events. The status view
reads these records; it does not issue new permissions.

A maintenance daemon converts repository events into jobs. It uses a SQLite
store in write-ahead-log mode for the events, jobs, run results, and resource
claims. Hooks, file scans, interruptions, and explicit commands feed the
store. Content hashes remove duplicate events; an allowlist restricts the jobs
that may run; a single-daemon guard prevents competing schedulers. During one
recorded snapshot, the status page reported the daemon as not running while
queued jobs remained visible. The queue persisted, but the service was inactive.

\subsection{What continuous mathematical work looks like}

A continuous goal resumes from the latest accepted research state: sources
read, routes eliminated, computations interpreted, and formal obligations
still open. Each run selects a useful next step, such as a lemma, a
counterexample, or a sharper reduction, and records what changed.

A representative internal trace makes this less abstract. It is an internal
agent rehearsal operated by the author, and this checkout carries no public
receipt for its counts. Over fifteen
successive turns, one problem-directed run recorded 313 visible progress
updates and 3,491 command events, with linked workers exploring separate
analytic, computational, and formal lanes. The movement was not a straight
line from prompt to proof. Numerical probes exposed structure; exact arithmetic
replaced promising samples; several attractive strengthenings were killed by
counterexamples; partial inequalities were narrowed to their valid domains;
Lean targets were attempted only after an ordinary mathematical kernel had
stabilised; and successful local results were propagated into the problem
frontier before the next search began. While long exact computations ran, the
controller advanced independent work instead of repeatedly polling them.

That trace also displays why the architecture keeps several evidence layers.
Its compact narrative reported no observed changed paths even though its own
turn summaries referred to landed commits. This is not evidence that nothing
changed, nor that the summaries were correct. It is evidence that a compressed
trace has an observation boundary. Repository state, command results, diffs,
formal builds, and commit objects must settle the question. Likewise, four
turns lacked complete closeout events and only sixteen of forty-two linked
session outcomes appeared in the compact projection. Completeness is therefore
an explicit field, not an impression created by fluent prose.

Each advance leaves an authority-bearing artefact and receipt. When the next
step needs an unavailable validator, a human decision, or a new idea, the run
preserves its unfinished obligation and a useful re-entry point.

\subsection{Coupled continuous goals: discovery and stewardship}
\papersectiontarget{systems-coupled-goals}

The public workflow separates proof search from maintaining the account of
what has been found. A discovery task investigates one mathematical question.
A stewardship task compares the result with the existing literature and
formal development, decides where it belongs, and updates the affected paper
and records. One agent may do these jobs in succession, or separate agents may
handle them. The separation makes time for questions that proof search alone
can neglect: whether the result is already known, whether several declarations
express one theorem, and whether the paper explains the hard step.

\begin{figure}[!t]
\centering
\begin{tikzpicture}[
  role/.style={paper stage, text width=3.05cm, minimum height=1.5cm},
  object/.style={role, fill=TealBlue!6, draw=TealBlue!60!black,
    text width=2.72cm},
  output/.style={role, fill=ForestGreen!6, draw=ForestGreen!55!black,
    text width=3.15cm},
  arr/.style={paper arrow},
  feedback/.style={paper feedback}
]
\node[role, fill=MidnightBlue!5, draw=MidnightBlue!55] (mine) at (0,0)
  {{\scriptsize\sffamily\bfseries DISCOVERY}\\[3pt]\textbf{Test the frontier}\\[2pt]
   \scriptsize compute, prove, or falsify};
\node[object, right=7mm of mine] (delta)
  {{\scriptsize\sffamily\bfseries HANDOFF}\\[3pt]\textbf{Recorded result}\\[2pt]
   \scriptsize theorem, no-go, computation};
\node[role, fill=BurntOrange!7, draw=BurntOrange!65!black,
  right=7mm of delta] (steward)
  {{\scriptsize\sffamily\bfseries STEWARDSHIP}\\[3pt]\textbf{Place the result}\\[2pt]
   \scriptsize appraise, reconcile, explain};
\node[output, below=9mm of delta] (surfaces)
  {{\scriptsize\sffamily\bfseries PUBLIC SURFACE}\\[3pt]
   \textbf{Current mathematical picture}\\[2pt]
   \scriptsize paper order, evidence, frontier};
\draw[arr] (mine) -- (delta);
\draw[arr] (delta) -- (steward);
\draw[arr] (steward) -- (surfaces);
\draw[feedback] (surfaces.west) to[out=180,in=-90]
  node[below left, font=\tiny\sffamily, align=center]
  {ranked next question, missing evidence,\\discriminating test} (mine.south);
\end{tikzpicture}
\caption{Search produces a result; a second pass checks its place in the existing
work and updates the account used to select the next question.}
\label{fig:coupled-goals}
\end{figure}

A stewardship pass can change the next research question. An apparently new
result may be a routine corollary; an overlooked theorem may already supply a
missing premise. The pass records the evidence, judges mathematical interest,
revises the explanation, and recommends further work. Proof status and
editorial prominence are recorded separately.

The \texttt{run-coupled-research-goals} skill invokes the search and update
jobs with a shared source revision. New results or corrections can trigger
another pass. If the repository has not changed, the jobs yield. Contributors
who need an introduction can first ask \texttt{explain-public-system} to read
the papers and explain one result at the requested level.

\section{The mathematical reasoning loop}\label{sec:mathloop}
\papersectiontarget{systems-mathloop}

Research begins with the mathematical question and the relevant literature.
The agent compares possible approaches, tests the promising ones, and
formalises steps once their statements are stable. The choice of approach is
an authored judgement: a counterexample that rules out an expensive line of
work may deserve more attention than many easily proved lemmas.

\subsection{Experiments are route selectors}

Computation serves three useful roles. It can find a counterexample, measure a
finite pattern, or test whether a proposed mechanism is plausible enough to
formalise. Each experiment records its inputs, code, finite domain, output, and
interpretation. The interpretation must state what the experiment cannot show.

A finite computation becomes formal evidence only through an explicit bridge.
For example, Lean may evaluate a finite proposition with \texttt{decide} and
check the resulting proof term. Even then, the conclusion remains finite. A
pattern observed for many inputs does not become an ``all inputs'' theorem, and
a successful numerical approximation does not become an exact equality. Failed
experiments are kept when they prune a natural route; otherwise later agents
pay to repeat the same mistake.

\paragraph{Failure modes are mathematical outputs.}
The system distinguishes three kinds of negative evidence. A failed agent
attempt says only that one search path did not close. A counterexample can
refute a conjecture on its exact domain. A Lean no-go theorem can rule out a
whole strategy class under explicit hypotheses. The last two may be as
valuable as a positive lemma: they prevent repeated work and reveal which new
ingredient an endpoint requires. This corpus therefore keeps deep
problem-specific reasoning surfaces, proof gaps, coefficient-only and
fixed-precision obstructions, and corrected statements beside successful
theorems. Every no-go keeps its scope visible; failure of one mechanism is not
failure of every possible proof.

\subsection{Following a change of hypothesis}
\papersectiontarget{systems-research-attempt}

The composite-dilation work gives a concrete example. A support here selects
the exponents in a Mersenne subseries. A researcher wants to
extend a multiplication identity from supports consisting of primes to more
general supports. The public query tool retrieves the existing statement,
its source module, and the claim that uses it. The researcher can then inspect
which use of primality must be replaced, rather than beginning with the whole
library. The relevant
\href{https://github.com/wcook04/plectis-erdos/blob/fd53947f32d9146629be163521d3d5c7cb87a1a8/Erdos249257/CompositeDilationDefect.lean\#L30}%
{formal source} counts the extra support divisors introduced by multiplication
as an explicit defect term, and proves that this term vanishes on prime
support. The extension has therefore become a precise question about that
term.

During an attempt, the agent uses a session record to keep observations,
conjectures, plans, and interpretations together. When it asks Lean to test a
candidate, the proof workbench saves the submitted source, computes its digest,
runs the probe, and appends the result. A later success record must cite an
accepted probe whose stored source still matches that digest. The agent can
revise its next candidate while leaving the earlier attempt available for
inspection. A reviewer follows the resulting argument and checks that the
explanation describes the statement actually tested. The
\href{https://github.com/wcook04/plectis-erdos/blob/fd53947f32d9146629be163521d3d5c7cb87a1a8/scripts/proof_workbench.py}%
{public proof workbench} implements this record-and-check step.

The useful memory of the failed extension is more specific than ``composite
support did not work''. Its
\href{https://github.com/wcook04/plectis-erdos/blob/fd53947f32d9146629be163521d3d5c7cb87a1a8/docs/semantic/lab/receipts.json}%
{failure record} names the lost hypothesis, the two supporting declarations,
and the repair: retain the defect and prove it vanishes or is controlled on
the chosen support. It also names the surviving mechanism and the condition
for reopening the route. The theory-laboratory builder joins that record to
the declaration and semantic graphs; its checker verifies that the cited
objects and neighbouring mechanisms exist. The next agent receives both a
reason to abandon the unmodified identity and a concrete direction to study.

An accepted change can leave several files out of date.
The declaration index reads the Lean source; source-coordinate records join
that index to the papers; the corpus descriptor summarises the resulting
state; and the publication entry packet reads the final claim and publication
records. The
\href{https://github.com/wcook04/plectis-erdos/blob/fd53947f32d9146629be163521d3d5c7cb87a1a8/scripts/refresh_projections.py}%
{projection refresher} orders these builders by their dependencies. Release
checks compare the stored projections with what their owners would regenerate.
These updates keep the paper, query results, and formal source consistent for
the next contributor.

\subsection{Checking the result}

Three checks answer different questions. The literature establishes the
reported status of the problem. Lean checks the exact formal statement and
proof. Mathematical review compares that statement with the intended question
and with prior work. The record preserves the source and result of each check.

When an argument comes from the literature, the citation remains attached to
the formalisation. When an experiment or failed proof is retained, its scope
and interpretation remain attached too. This lets later researchers examine
why an approach was pursued and what the attempt established.

\subsection{From local progress to reusable mathematics}

A lemma proved for one Erd\H{o}s problem may be useful elsewhere. Preparing it
for a shared library requires another pass: determine which hypotheses the
proof uses, search for an existing theorem, and state the result in a form
that other developments can use. The general proof must be checked, and the
original result should follow as a clear special case.

The repository can preserve that starting point and help prepare a submission.
Mathlib's contributors and reviewers decide whether it belongs in the library
\cite{mathlib}. Credit should distinguish the original observation from the
generalisation, formalisation, and library work. The
\papersectionlink{open-source-mathematics-strategy.pdf}%
{strategy-local-to-general}{strategy paper} describes this proposed route.

\subsection{Problem-sized Lean worlds and bounded theorem neighbourhoods}

The declaration atlas locates formal statements. Dependency maps show which
declarations a proof uses, while authored relations explain their mathematical
role. The claim record links selected results to their public descriptions.
A researcher uses these views together when deciding which source to inspect.

The scale is already problem-sized. At the August 2026 snapshot the attached
public corpus described 1,024 Lean modules and 153,396 declarations. The much
deeper private record for Problem 1041 contained 177 exact results, seven open
producers, 72 negative results, 134 Lean modules, and 335 executable experiment
programs. Its public research-corpus manifest alone recorded 685 files and 35
selected strongest results. These are dated inventory facts, not a throughput
benchmark; generated certificate families also make raw theorem counts a poor
measure of mathematical value. The private figures come from the author's own
workbench, and this public checkout carries no receipt against which a reader
can check them.

At that snapshot, the canonical public 1041 library had two
integrated Lean modules, while its separately governed public research export
has 125 Lean sources and its private world is larger again. The public route
memory still records no reviewed 1041 claim family. Across the whole semantic
corpus, structural linkage is broad, but only 26 statement nodes and eight
relations carry digest-bound semantic-review receipts in this release, as
reported by \texttt{scripts/semantic\_review.py --check}. These
differences are intentional status boundaries: a discoverable declaration is
not thereby understood, integrated, reviewed, or publishable.

An agent does not receive this world as one prompt. First the problem cockpit
fixes the endpoint, status, claim ceiling, and canonical frontier. It keeps open
producers distinct from the target so a promising route cannot silently replace
the problem. Next a bounded neighbourhood selects the strongest relevant
results, exact source coordinates, imports, consumers, alternative formulations,
experiments, counterexamples, no-gos, and literature. It labels every item by
evidence class and emits an omission receipt with an expansion route. One 1041
packet exposed only four of 134 Lean modules, four of 335 experiments, and four
of 55 no-go notes; designed omission, not exhaustive loading, made the context
usable. When an exact path is claimed, a local connection card can put its
prerequisites, sibling mechanisms, consumers, falsifiers, and validation target
before broader retrieval.

\subsection{Comprehension before the mathematics}
\papersectiontarget{systems-comprehension}

The private working-memory compiler separates orientation from focused work.
An overview shows the available results and open questions so that an agent
can choose a target. It does not construct candidate proof routes. A focused
request names a claim, declaration, or premise; if no target can be resolved,
it is reported as unanchored and returns no proof branches.

Each attached corpus retains its own index; nothing is copied into a central
index. Compact descriptors and content fingerprints identify the sources.
The compiler opens a larger atlas when the bounded query does not supply
enough context. The focused packet starts with the target, established
results, and missing premises. Its workspace groups premises into candidate
arguments and records tests or counterexamples that could reject them.

A freshness check compares the recorded projection and descriptor fingerprints
with their current values and requests recompilation when they differ. An
exactly named Lean definition can also be newer than its index. In that case,
the compiler returns the current definition and use-site coordinates, and
defers broader graph-based routes until the index is refreshed.

When the selected material exceeds the context budget, the compiler first
shortens summaries and repeated branches while preserving source identities
and retrieval routes for omitted details. Only if the resulting packet still
does not fit does it report the requested and estimated size and exit
non-zero. These checks make the selected context inspectable. Whether it
leads to better mathematics requires the comparison proposed in
Section~\ref{sec:related}.

\subsection{Consequence propagation}

After a declaration changes, the agent checks the files that may depend on it:
imports, later proofs, experiments, claim records, and papers. The dependency
map identifies candidates; a semantic second pass determines which actually
need revision. Each is updated, checked as unchanged, deferred with a reason,
or excluded from the change. No projection may bulk-strengthen a family of
claims.

The public \texttt{propagate-research-consequences} skill gives instructions
for this pass. Work from an older clone is checked against its original
revision and again after integration with current main. The original
contribution and the integration changes retain separate attribution.

A recurring process failure may also justify changing a tool or instruction.
The proposed repair records the local failure, other affected cases, and a
check of the repair. The implementation supports this form of feedback;
it has not measured how often useful lessons are successfully retained.

\section{The public Lean repository}\label{sec:public}
\papersectiontarget{systems-public}

The public checkout begins at a deliberate boundary. It contains every file
needed to inspect its claims and replay its checks. It does not call back into
the private workbench, and no public theorem depends on an unpublished private
lemma. The larger workflow is provenance: it explains production, not truth.

The repository has two Lean roots. The reviewed root contains the established
249/257 publication lane. The problem-owned expansion root contains work on six
additional Erd\H{o}s problems and unpromoted lanes for 249 and 257. Both roots
are checked by the same Lean kernel, but kernel acceptance does not automatically
promote a declaration into a reviewed public claim. Promotion is a separate
change to the claim record and exposition.

The main public sources have deliberately separate jobs:

\begin{table}[H]
\centering\small
\begin{tabular}{@{}>{\raggedright\arraybackslash}p{0.23\linewidth}
                    >{\raggedright\arraybackslash}p{0.31\linewidth}
                    >{\raggedright\arraybackslash}p{0.34\linewidth}@{}}
\toprule
\papertablehead
Surface & Authority & It does not establish \\
\midrule
Lean source and pinned toolchain & Exact formal statements and proofs accepted
by the kernel. & Intended meaning, novelty, significance, or faithful prose. \\
\addlinespace
Reviewed claim record & Approved wording, status, evidence, bounded domain,
and adjacent open statement for selected claims. & Correctness or completeness
of the human review. \\
\addlinespace
Papers and guides & Human explanation and reading order within the recorded
claim ceiling. & New formal authority. \\
\addlinespace
Generated maps and query tools & Bounded navigation across declarations,
problems, graphs, papers, and claims. & Proof or permission to strengthen a
claim. \\
\addlinespace
Release checks and continuous integration & That configured identities,
relationships, generated files, licences, and negative tests pass on a named
revision. & Understanding of unrestricted prose or independent mathematical
approval. \\
\bottomrule
\end{tabular}
\caption{The public authority split.}
\label{tab:authority}
\end{table}

The companion public Plectis repository has a different role. It publishes
runnable, bounded mechanism slices and receipts from the wider system. The Lean
repository publishes a mathematical corpus. Neither public repository inherits
authority from the other, and neither is a public mirror of the private root.

\section{Comparator, Palomar, and publication}\label{sec:assurance}

\subsection{Comparator: an exact-statement firewall}

Comparator protects a narrow but important boundary. Selected propositions are
declared again in a challenge module without their proofs. A solution module
must provide terms of those exact types. The configuration pins the permitted
axioms, the modules, and a runtime receipt. A named altered statement must fail,
which checks that the harness has not become vacuous.

Comparator therefore answers: ``Does the proof-bearing corpus still implement
this separately stated Lean interface under this axiom budget?'' It does not
answer whether the interface is the right translation of an informal problem,
whether the theorem is new or important, or whether an independent human has
reviewed the proof. ``Comparator-checked'' is accurate; ``independently
verified'' is not.

\subsection{Palomar: selecting what deserves review}

Comparator's roster supplies evidence for a review portfolio. To prepare that
portfolio, maintainers compare result families in a local selection record.
Each assessment explains the consequence, hard mechanism, evidence, and open
remainder, then records why the family is selected, represented by another
result, deferred, or set aside. The local Palomar qualification checker checks
that this record covers its declared candidates and that the selected packet
satisfies the submission requirements. The mathematical ranking is authored;
the checker validates its supporting record.

This separation matters under proof abundance. Counting theorems rewards
generated volume and routine closure. Selection instead asks which result changes
the mathematical picture, which conditional route is closest to an endpoint,
which obstruction prevents wasted work, and which explanation will help an
expert assess the claim. The ranking may guide attention; it cannot alter proof
status.

\subsection{Publication and propagation}

After formal proof and statement reconciliation, a result may enter an authored
paper, a claim record, a Comparator packet, and a Palomar review unit. Each is a
separate projection with a separate ceiling. A public release is made only
after the Lean build and the release-surface checks pass, and publication
remains a human action.

The return path is equally important. A proof failure may improve a
formalisation heuristic. A counterexample may close a family of tempting
routes. A reviewer objection may require a narrower public claim. A release
drift may create a stronger check. Those lessons propagate to the smallest
durable owner: a theorem, experiment record, skill, standard, route, or claim
boundary. They do not rewrite raw intent, generated views, or mathematical
status by implication.

This places the repository inside a longer mathematical lifecycle:
candidate, proof, exact checking, exposition, independent assurance,
publication, expert digestion, and eventual community canonicalisation
\cite{taoai}. The system directly supports the early and middle stages. It can
prepare later stages but cannot award them to itself.

\subsection{What the papers should explain}

Tao's essay follows proof generation, verification, exposition, publication and
community digestion through to the adoption of useful mathematics
\cite{taoai}. His discussion of exposition is especially relevant to a
repository largely written by agents. Fluent prose can spend pages on routine
steps while hiding the step that took real work. It can also erase the
\emph{natural friction} that tells a reader where to slow down.

The problem papers should therefore explain why an approach was chosen, where
it becomes difficult, and what an unsuccessful attempt teaches. The failure
records supply evidence for that explanation: a lost hypothesis, an exact
counterexample, or a proved obstruction. Tao's later discussions make the case
for preserving such insight while a problem is still open
\cite{taolearning2026,taoexploration2026}.

Writing can reveal errors in the research record. An unexplained implication
may conceal a missing premise; a confusing paragraph may combine distinct
results. The author must return to the statements and repair the account.

The writing skill also asks the agent to compare the manuscript's selection
with the actual source tree. A useful theorem can be missing from the paper
or from the query that supplied its outline. The agent reads omitted results
and their uses before deciding whether the account should change. When several
proved ingredients appear to settle a step still described as open, it writes
out their composition and checks the assumptions at each connection. If that
argument needs further formalisation, the missing step becomes a proof task.
The paper can explain the ordinary argument while identifying precisely which
part has passed Lean.

A corrected account must reach the reader. The agent updates the existing
result records and reading order, rebuilds the affected views, and inspects the
actual query output. A registry edit alone can leave a result absent from the
front door. During concurrent work, publication views are generated from a
committed snapshot so that another agent's unfinished proof cannot enter the
published evidence accidentally.

This paper was itself produced within the system. Its evaluation is
author-operated.

\section{One complete boundary: finite is not unbounded}\label{sec:example}

Erd\H{o}s Problem 249 asks whether
\[
S=\sum_{n\ge1}\frac{\varphi(n)}{2^n}
\]
is irrational, where \(\varphi(n)\) is Euler's totient function
\cite{erdosgraham}. The development reduces irrationality to a family of exact
non-integrality certificates. Along the diagonal
\(H_t=\operatorname{lcm}(1,\ldots,t)\), write \(\mathrm{Cert}(t)\) for the
existence of the required finite certificate.

\href{https://github.com/wcook04/plectis-erdos/blob/fd53947f32d9146629be163521d3d5c7cb87a1a8/ErdosProblems/Skip/LadderT67.lean\#L71264}%
{Lean checks the finite theorem}
\[
\forall t\le82,\quad \mathrm{Cert}(t).
\]
The endpoint needs an unbounded supply:
\[
\forall T,\quad \exists t>T,\quad \mathrm{Cert}(t).
\]
The development proves that the unbounded statement is
\href{https://github.com/wcook04/plectis-erdos/blob/fd53947f32d9146629be163521d3d5c7cb87a1a8/Erdos249257/LcmConeFlatness.lean\#L426}%
{equivalent to the irrationality claim}. It does not prove the missing implication from the finite
range to the unbounded statement. No matter how large a fixed checked bound is,
a larger cutoff exists.

This example passes through every layer. Computation constructs finite data.
Lean checks the certificate predicate and the finite theorem. The formal graph
records its dependencies. The semantic graph identifies its role. The claim
record states the finite range and names the unbounded requirement as open.
The paper explains why the quantifiers differ. Comparator checks selected exact
interfaces. Palomar may rank the result family for review. The release checker
requires the limitation to remain visible.

The last check was added because the boundary once escaped. A historical README
edit changed a clause saying that the finite cases did \emph{not} supply the
open requirement into one saying that they completed it. Lean was untouched,
and the release checker passed because that prose relationship had not been
registered. After the relationship was added to the claim record, a deliberately
false copy was rejected. This is evidence for one failure and repair. It is not
a detection rate, a proof that every claim-bearing sentence is registered, or
an evaluation of mathematical review quality.

The historical study was narrower than a benchmark. It was a deterministic
mutation test: a known false edit was applied to a pinned tree, and the
configured release checks either rejected it or passed. It reports the reach of
those checks on those edits, and it reports nothing about how a human reader
comprehends the corpus. Its protocol, matrix, and limits are recorded in
\path{docs/publication_evidence.json}. In that study, nine of the
ten edits were rejected. One escaped because the relationship had not been registered. The
original run logs were not retained. The edits were authored by the checker's
author. After the relationship was registered, only the escaped edit was
reconstructed against the repaired checklist. The other nine edits were not
rerun against the extended checklist. The post-repair witness accepts the
current README and
rejects a test copy containing the false clause.

The same architecture handles other shapes of progress. Problem 257 separates
\href{https://github.com/wcook04/plectis-erdos/blob/fd53947f32d9146629be163521d3d5c7cb87a1a8/Erdos249257/CertificateKernel.lean\#L8328}%
{a theorem for full-support representations} from a stronger arbitrary-support
statement that remains open. Problem 68 has \href{https://github.com/wcook04/plectis-erdos/blob/fd53947f32d9146629be163521d3d5c7cb87a1a8/ErdosProblems/Erdos68/FactorialZeroPlateau.lean\#L971}%
{an exact carry-based equivalence}
but no theorem producing the required carries. Problem 251 has \href{https://github.com/wcook04/plectis-erdos/blob/fd53947f32d9146629be163521d3d5c7cb87a1a8/ErdosProblems/Erdos251/PrimeGapDyadicTail.lean\#L435}%
{an exact series reformulation} rather than an irrationality proof, and Problem 243 records a
\href{https://github.com/wcook04/plectis-erdos/blob/fd53947f32d9146629be163521d3d5c7cb87a1a8/ErdosProblems/Erdos243/SparseResetRecovery.lean\#L510}%
{conditional recovery theorem} whose premises remain part of the public claim.
A counterexample or corrected statement, as in the 1041 lane, is also a
first-class result. The publication rule is identical in every case: preserve
the quantifiers, premises, and nearest stronger open statement.

\section{Inspection routes}\label{sec:routes}

A reader can inspect the public system through short, question-shaped routes.
The labels below link to repository paths; the paper does not print the full
URLs.

\begin{table}[H]
\centering\small
\begin{tabular}{@{}>{\raggedright\arraybackslash}p{0.34\linewidth}
                    >{\raggedright\arraybackslash}p{0.54\linewidth}@{}}
\toprule
\papertablehead
Question & Start here \\
\midrule
How does the public repository fit together? &
\repolink{ARCHITECTURE.md}{ARCHITECTURE.md} \\
What may the project say, and what remains open? &
\repolink{docs/claims.json}{docs/claims.json} and
\repolink{docs/methodology.json}{docs/methodology.json} \\
Where is the checked mathematics? &
\repolink{Erdos249257.lean}{Erdos249257.lean} and
\repolink{ErdosProblems.lean}{ErdosProblems.lean} \\
How can I navigate without reading the whole corpus? &
\repolink{docs/ORIENTATION.md}{docs/ORIENTATION.md} and
\repolink{scripts/query_corpus.py}{scripts/query\_corpus.py} \\
What exactly does Comparator check? &
\repolink{docs/EXTERNAL_VERIFICATION.md}{docs/EXTERNAL\_VERIFICATION.md} and
\repolink{verification/comparator.json}{verification/comparator.json} \\
What does Palomar select and qualify? &
\repolink{docs/PALOMAR_QUALIFICATION.md}{docs/PALOMAR\_QUALIFICATION.md} and
\repolink{docs/PALOMAR_RESULT_SHOWCASE.json}{docs/PALOMAR\_RESULT\_SHOWCASE.json} \\
Which papers exist and what question does each answer? &
\repolink{docs/papers/README.md}{docs/papers/README.md} \\
Which checks gate a release? &
\repolink{scripts/check_release.py}{scripts/check\_release.py} and
\repolink{.github/workflows/lean.yml}{.github/workflows/lean.yml} \\
How can work return with public credit? &
\repolink{CONTRIBUTING.md}{CONTRIBUTING.md} and
\repolink{docs/research-commons/CONTRIBUTIONS.md}{research-commons/CONTRIBUTIONS.md} \\
\bottomrule
\end{tabular}
\caption{Public files for inspecting the source, checks, and contribution record.}
\label{tab:routes}
\end{table}

\section{What can be trusted}\label{sec:trust}
\papersectiontarget{systems-trust}

The architecture supports a chain of narrow conclusions:

\begin{enumerate}
\item a recorded experiment ran on its stated finite inputs;
\item a pinned Lean kernel accepted its stated formal source;
\item a maintainer approved a selected interpretation and public boundary;
\item Comparator matched selected proof-bearing declarations to separately
      stated interfaces and an axiom budget;
\item Palomar's local checks validated the structure of a selected review
      packet;
\item the release program found every configured public relationship intact on
      the named revision.
\end{enumerate}

The checks have a common limitation: agreement between files is weaker than
independent examination of their content. A coordinated error in source,
claim record, and prose can pass a structural comparison. Unregistered prose
can overclaim, and literature status can change. The system does not require
a second independent mathematician to approve every result.

This is also a limitation of assurance cases: recording a relation between
evidence and a claim does not establish that the evidence is sufficient
\cite{gsn}. The links and negative tests make selected claims easier to
inspect, while mathematical judgement remains necessary.

\section{Scaling from one clone to a search network}\label{sec:scaling}
\papersectiontarget{systems-scaling}

\subsection{Clone, attach compute, and mine bounded results}

A contributor can use a different agent runner with the public source and
checks. Several workers can investigate separate statements or experiments,
then return their results for integration. The repository calls this
\emph{result mining}. Each return names the result, evidence, starting
revision, and unresolved question.

As the number of workers grows, someone must still compare their returns with
the literature and existing corpus, update the papers, and choose the next
questions. This is the stewardship work described earlier. A large run may
produce no new result; a short counterexample may change the research plan.
The number of workers does not determine the amount of useful work to publish.

The controller limits concurrent work according to file conflicts and machine
capacity. Reading and small calculations can proceed while another task holds
the Lean validation slot.

Lean validation uses semantic single-flight coordination. A request is identified
by the reachable source, logical targets, toolchain, dependency lock, and
authority mode. Its command string and working directory do not determine
request identity.
Equivalent requests share one owner and may reuse a completed receipt; changed
inputs require a fresh build. Distinct requests that would still compete for
the same host-wide Mathlib resource are serialized by a shared conflict key.
A distinct request that cannot acquire the slot ends with a resource-busy
deferral before its validator starts. The agent can continue independent work,
then resubmit that request after the resource is free. Waiting on the deferral
receipt does not start a build. An equivalent request already running can be
joined instead.

These are four separate scaling limits: the number of research attempts,
safe concurrent editing, proof-check throughput, and expert review. Increasing
one can leave another as the bottleneck.

\paragraph{Returning work and recording credit.}
A contributor can fork or clone the repository and submit a pull request or
research-progress issue. The return records the result, starting commit,
evidence, collaborators, and tools used. A maintainer supplies an explicit
review decision and the commit containing the accepted work. The acceptance
program checks the submission and commit ancestry, then copies the return into
an accepted receipt while preserving its result, evidence, and attribution.
Only an accepted receipt enters the generated contribution views. Their builder
reads receipts committed in the checked-out history and rejects a receipt
whose working copy differs from that committed source.
The published credit therefore identifies the accepted work and the review
decision that admitted it. Acceptance, mathematical claim status, and release
inclusion remain separate decisions. Later
corrections preserve the earlier record and its attribution.

The starting commit lets a maintainer reproduce the original contribution
before adapting it to current main. A substantive conflict resolution receives
separate integration credit. The
\papersectionlink{open-source-mathematics-strategy.pdf}%
{strategy-protocol}{contribution protocol} and
\papersectionlink{open-source-mathematics-strategy.pdf}%
{strategy-credit}{credit policy} give the details. The
\papersectionlink{cold-clone-to-proof-receipt.pdf}%
{cold-clone-problem}{cold-clone study} examines the navigation used to begin this work.

The public clone supplies the corpus and checks, but the private orchestration
system is not a turnkey public service. Large multi-provider deployments remain
a design target rather than a reported benchmark. No outside contributor had
completed this path by 31 August 2026.

\subsection{Using records of unsuccessful approaches}

A recorded obstruction should identify the method and hypotheses it excludes,
along with variants that remain possible. Linking these records to theorems
and experiments may help an agent avoid repeating a failed approach or spot a
related obstruction in another problem.

Graph-conditioned models could be tested on these records. For example,
training data might pair a tempting argument with a counterexample or theorem
that explains its failure. These are proposals for future systems. A mistaken
or over-broad relation could instead hide a viable method, so an evaluation
would need reviewed source records and held-out tasks.

\subsection{An experimental agenda}

Three questions must be kept apart. The first asks whether the runtime at a
past revision exposes the exact remaining obligation for a stated target and
accepts an authored candidate that avoids the later declaration and module
names. The second asks whether the semantic layers present at that revision
help a reader find existing mathematics and interpret it correctly. The third
asks whether the derived layers as they stood in the past shorten the work of
producing a later result. The harness implemented here answers the first. It
supplies its own authored candidate, names the candidate declaration inside its
request, runs only the derived-layer arm, and observes no reader and no agent,
so it leaves the second and the third open. Its clone carries the past source
with the dependencies of the current tree, so its receipt records the pins that
revision declared alongside the pins the build used. Later Git objects, a
shared object store, generated dossiers and receipts, network access, context
supplied in the prompt, and a model's familiarity with public mathematics can
each reveal the target. \path{docs/methodology.json} enumerates these channels
and records which the harness controls. The channels it cannot control are
limitations of any result produced this way.

The public benchmark builder prepares source snapshots from before a target
declaration was introduced. It checks that the declaration is absent and
filters added graph and mechanism records against the declarations available
at that revision. The snapshot omits Git history, since a linked worktree would
let an agent recover the later proof from the shared object store. The answer
key and introducing-commit metadata remain outside the packet. These are
prepared inputs for an evaluation; the runner must still isolate filesystem
and network access, and the selected material needs review for answer hints.

The next evaluation should ask more than how many theorems additional compute
produces. It should compare graph-aware and graph-blind agents on frontier
selection, repeated-dead-end rate, time to a useful obstruction, premise
discovery, proof completion, and calibration of public claims. It should test
transfer between unrelated mathematical domains as well as between neighbouring
Erd\H{o}s problems. Independent teams should replay result packets and review
whether no-go edges are scoped correctly. The central scaling question is
whether an increasingly rich map of what works and what cannot work makes
future reasoning more efficient and more original, or merely makes the system
more confident about the territory it already knows.

\section{Relation to other approaches}\label{sec:related}
\small

\paragraph{Theorem-proving agents.}
This architecture is not a new proof-search algorithm. LeanDojo and Pantograph
provide programmatic Lean environments and proof-state interaction
\cite{leandojo,pantograph}. OpenProver already combines a planner, parallel
workers, independent verifiers, compact working memory, a larger repository,
and Lean feedback \cite{openprover}. Agent Hunt already studies concurrent
formalisation with locks, bounties, guarded ownership, and collaborative
agents \cite{agenthunt}. LeanMarathon uses an evolving Lean blueprint as both proof graph and shared
record, with contract-scoped agents, adversarial target review, and parallel
CI-gated proof work \cite{leanmarathon}. It also stores explanatory prose beside
Lean types and checks agreement between cited and elaborated proof dependencies.
This overlaps both the durability and checked-exposition concerns here.
DreamProver learns a compact reusable lemma library
through wake--sleep cycles \cite{dreamprover}. These are baselines for proof
interaction, parallel search, and learned proof memory. The present work follows the result into the public research record, including
its explanation, open obligations, and contribution history.

\paragraph{Graphs and checked exposition.}
Proof blueprints connect informal proof plans to named Lean declarations.
\texttt{leanblueprint} records the declarations that realise a statement, the
statement and proof dependency edges its author declares, and a
formalisation-readiness state for every node, and its declaration check
confirms that each named declaration exists in the project or a dependency
\cite{leanblueprint}. LeanArchitect infers formal dependencies and
unfinished-proof status and exports synchronised blueprint material
\cite{leanarchitect}. Those systems already supply the declaration-linking
layer described above. The graph layers here have a related navigational role.
What this repository adds is the maintained relationship among public claims,
several evidence classes held in one record, obstructions that are themselves
proved, the generated projections a reader navigates, and a research task that
stays open after the document is published; the public-claim
record begins after a result has been selected and asks which wording and open
boundary were reviewed.

Semantic-audit systems address a neighbouring problem. Lean Atlas narrows the
declarations a person must inspect for chosen theorem statements, conditional
on the semantic correctness of the returned set and trusted base
\cite{leanatlas}. EconCSLib uses models to translate and compare formal and
informal statements, with saved human judgements for a subset \cite{econcs}.
The present architecture uses models throughout production but assigns none of
them final semantic authority. Their output becomes evidence or a candidate
until a source-specific gate accepts it.

Checked-document and traceability systems make heterogeneous relationships
explicit. Isabelle/DOF places formal and informal material in a typed checked
document \cite{isadof}; requirements traceability follows commitments across
development and revision \cite{gotel}. This repository keeps prose unrestricted
and records selected public boundaries separately. That choice makes adoption
simple, but leaves unregistered prose outside the checker.

\paragraph{Auditable scientific agents.}
HEP makes hypotheses, evidence, belief updates, lineage, and resolution states
explicit in an append-only registry \cite{hep}. Symposium proposes immutable
community publication histories with attributable artefacts, declared
evidence, assumptions, and purpose-sensitive arguments \cite{symposium}.
EurekAgent treats permissions, artefacts, budgets, and human supervision as
first-class parts of an agent environment \cite{eurekagent}. Persistent records
and environment design are therefore not unique to this system. The narrower
distinction here is between authorities that must not be collapsed:
experimental evidence, kernel acceptance, reviewed interpretation,
exact-statement assurance, editorial selection, and public release.

Mutation testing asks whether a test distinguishes a seeded fault from the
original program \cite{demillo,jiaharman}. The deliberately false boundary in
the worked example follows that idea. Because the examples were selected by
the system's author and were not run as a controlled external evaluation, they
show that particular checks can fail; they do not measure overall adequacy,
and a deterministic mutation outcome is not a measurement of human
comprehension.

\paragraph{Contribution and limits.}
The contribution is the implemented connection between ongoing research and
its public mathematical record: formal statements, interpreted results, scoped
failures, contribution credit, and checked publication relationships.
LeanMarathon already addresses durable agent coordination and statement
fidelity; \texttt{leanblueprint} already supplies statement-to-declaration
links, author-declared dependency edges, and formalisation-readiness state;
HEP already records hypotheses and refutations. These overlaps make
component-level comparison more informative than a claim to cover the most
stages. The evidence reported here is a design, internal agent rehearsals
operated by the author, and deterministic mutation tests over one evolving
corpus. It is not a controlled comparison of research productivity, and it is
not a measurement of reader comprehension.

One further piece of evidence exists and carries the same ceiling. On
6~September 2026 eight bounded agent readers, one per task and access arm,
read the frozen release candidate under four held-out tasks (an
ordinary-proof boundary, an inherited statement with a local proof, a finite
bound beside an unbounded theorem, and a target-equivalent reformulation).
One arm saw only the Lean tree, the manuscripts and the root prose; the other
arm also saw the claim register, the semantic corpus and the query tool.
Returns were adjudicated blind to arm against the cited declarations and the
manuscripts. All eight reached the strongest supported statement with its
hypotheses, none overstated a formal component, and the two arms did not
separate at this sample size. The readers found four defects in the corpus,
each repaired at its owner before this text was assembled, and the receipt at
\texttt{docs/measurements/fresh\_context\_use\_experiment\_20260906.json}
records the design, the scores and the defects. This is internal agent
evidence over four tasks; it says nothing about human readers.

\paragraph{Scaling beyond this corpus.}
The design is not tied to Erd\H{o}s problems. A new domain needs stable result
identities, source and status records, a formal checker where one exists, and
an explicit public-claim boundary; it need not use the same mathematics or
even Lean. Larger deployments should federate independently owned corpora
rather than build one authority database, separate producer and reviewer
teams, benchmark each boundary independently, and preserve immutable result
generations. Proof-search layers can adopt distributed task markets,
persistent lemma learning, and richer human steering, while publication layers
can add external replay, signed review receipts, and cross-project claim
links. None of these extensions should allow learned process memory or a graph
edge to change mathematical truth by itself.
\normalsize

\section{Conclusion}\label{sec:conclusion}

Plectis connects an ongoing research process to a public record that another
contributor can continue. The working example shows the intended unit of
reuse: a failed extension leaves a correction term, its formal statement,
and a precise next question. The surrounding machinery records attempts,
checks statements, and refreshes the linked explanations and source views.

The present evidence is a working implementation and author-operated cases.
The next evaluation is whether outside contributors can recover the relevant
mathematics, return reproducible work, and reduce the effort required for the
following attempt. That test should measure understanding and reuse as well
as proof completion.

\appendix
\small
\section{Reproducibility}\label{app:repro}

\paragraph{Public first contact.}
A fresh public clone can reproduce the control card and structural checks:
\begin{verbatim}
python3 scripts/proof_cockpit.py --format card
python3 scripts/proof_cockpit.py --check
\end{verbatim}
The first reads committed public metadata; the second checks the claim
registry, cold-clone contract, and orientation freshness.  Neither checks a
proof.  Formal authority begins with the pinned Lean build named by the card.

\paragraph{Dated navigation counts.}
At the semantic review's snapshot, taken before the eight-problem
consolidation of 2026-08-02 enlarged the corpus to 153,238 declarations, all
151,761 then-live declarations were inventoried and routed, and all
\SemanticAuthoredTheoremLike\ author-written theorem-like
declarations had an exact node link.  Of those,
\SemanticAuthoredInterpreted\ (\SemanticAuthoredPercent\%) participated in
authored mathematical interpretations: \SemanticDirectEvidence\ as exact
proposition evidence and \SemanticContextual\ as bounded certificate- or
module-family context.  The remaining \SemanticStructuralOnly\ were linked only
through exact source-module and normalised-signature families, not authored
mathematical paraphrases.  Every declaration selected for a public claim had
an authored route.  The command
\texttt{python3 scripts/query\_semantic.py coverage} derives these volatile
navigation counts and checks their references.  They do not measure semantic
review quality or public-claim completeness.

The paper inventory, \path{docs/publication_contract.json}, records source
and PDF cryptographic hashes and validation commands. The worked-example evidence is recorded in
\path{docs/publication_evidence.json}. The reconstruction file
\path{experiments/publication_mutations.json} specifies the deliberately false
edits and their limitations. Comparator's public configuration and receipt are
named by \path{verification/comparator.json}; Palomar's local readiness and
ranking surfaces are named by the corresponding files under \path{docs/}.
These artefacts identify what was checked. They do not interpret unrestricted
prose or confer external acceptance.

The private system is described here at the architectural level because it is
production provenance, not a dependency of the public result. Private paths,
operator material, unreleased work, and private ledgers are neither required
nor granted authority by this paper. The public checkout remains the inspection
and replay boundary.

{\tiny
\begin{multicols}{2}
\begin{thebibliography}{20}

\bibitem{leanmarathon}
Y. Zhang, Y. Sun, T. Suzuki, J. D. Lee, and F. Liu,
\emph{LeanMarathon: Toward Reliable AI Co-Mathematicians through Long-Horizon
Lean Autoformalization}, 2026,
\href{https://arxiv.org/abs/2606.05400}{arXiv:2606.05400}, Sections 2--4.

\setlength{\itemsep}{0pt}
\setlength{\parsep}{0pt}
\bibitem{lean4} L.~de Moura and S.~Ullrich, \emph{The Lean 4 Theorem Prover
and Programming Language}, in \emph{Automated Deduction---CADE 28}, Lecture
Notes in Computer Science 12699, 2021, pp.~625--635,
\href{https://doi.org/10.1007/978-3-030-79876-5_37}{DOI}.
\bibitem{erdosgraham} P.~Erd\H{o}s and R.~L.~Graham, \emph{Old and New
Problems and Results in Combinatorial Number Theory}, Monographies de
L'Enseignement Math\'ematique 28, 1980, p.~61.
\bibitem{leanblueprint} P.~Massot, \emph{leanblueprint}, plasTeX plugin for
Lean formalisation blueprints, 2020,
\href{https://github.com/PatrickMassot/leanblueprint}{software repository}.
\bibitem{leanarchitect} T.~Zhu, P.~Monticone, S.~Welleck, and J.~Avigad,
\emph{LeanArchitect: Automating Blueprint Generation for Humans and AI},
in \emph{17th International Conference on Interactive Theorem Proving},
LIPIcs 382, 2026, pp.~25:1--25:16.
\bibitem{leanatlas} B.~Yanahama and A.~Sannai, \emph{Lean Atlas: An
Integrated Proof Environment for Scalable Human--AI Collaborative
Formalization}, 2026, \href{https://doi.org/10.48550/arXiv.2604.16347}{arXiv}.
\bibitem{econcs} N.~Garg, \emph{EconCSLib: AI-Assisted Lean Formalization for
Economics \& Computation Research}, 2026,
\href{https://doi.org/10.48550/arXiv.2606.13306}{arXiv}.
\bibitem{isadof} A.~D.~Brucker and B.~Wolff, \emph{Isabelle/DOF: Design and
Implementation}, in \emph{Software Engineering and Formal Methods}, 2019,
pp.~275--293.
\bibitem{gotel} O.~C.~Z.~Gotel and A.~C.~W.~Finkelstein, \emph{An analysis
of the requirements traceability problem}, Proc. First IEEE International
Conference on Requirements Engineering, 1994, pp.~94--101.
\bibitem{demillo} R.~A.~DeMillo, R.~J.~Lipton, and F.~G.~Sayward, \emph{Hints
on test data selection}, IEEE Computer 11(4), 1978, pp.~34--41.
\bibitem{jiaharman} Y.~Jia and M.~Harman, \emph{An analysis and survey of the
development of mutation testing}, IEEE Transactions on Software Engineering
37(5), 2011, pp.~649--678.
\bibitem{gsn} SCSC Assurance Case Working Group, \emph{Goal Structuring
Notation Community Standard, Version 3}, SCSC-141C, May 2021.
\bibitem{taoai} T.~Tao, \emph{Mathematics in the age of AI}, preprint, 2026,
\href{https://doi.org/10.48550/arXiv.2608.16753}{arXiv}.
\bibitem{taolearning2026}
T.~Tao, discussion of learning from studying a mathematical problem,
Mastodon thread, 3 September 2026,
\href{https://mathstodon.xyz/@tao/117208617602946453}{opening post},
accessed 5 September 2026.
\bibitem{taoexploration2026}
T.~Tao, discussion of exploration and insight in the Navier--Stokes
regularity problem, Mastodon thread, 3 September 2026,
\href{https://mathstodon.xyz/@tao/117207849921390904}{opening post};
\href{https://mathstodon.xyz/@tao/117207855800042681}{part 5 on preserving exploration},
accessed 5 September 2026.

\bibitem{mathlib} Lean community, \emph{Contributing to mathlib},
\href{https://leanprover-community.github.io/contribute/index.html}%
{contributor guide}, accessed August 2026.
\bibitem{leandojo} K.~Yang et al., \emph{LeanDojo: Theorem Proving with
Retrieval-Augmented Language Models}, NeurIPS 2023.
\bibitem{pantograph} J.~Storrs et al., \emph{Pantograph: A Machine-to-Machine
Interaction Interface for Advanced Theorem Proving, High Level Reasoning, and
Data Extraction in Lean 4}, 2024,
\href{https://doi.org/10.48550/arXiv.2410.16429}{arXiv}.
\bibitem{openprover} M.~Kripner and M.~Straka, \emph{OpenProver: Agentic and
Interactive Theorem Proving with Lean 4}, 2026,
\href{https://doi.org/10.48550/arXiv.2607.09217}{arXiv}.
\bibitem{agenthunt} C.~E.~Brown, C.~Kaliszyk, and J.~Urban, \emph{Agent Hunt:
Bounty Based Collaborative Autoformalization With LLM Agents}, 2026,
\href{https://doi.org/10.48550/arXiv.2603.06737}{arXiv}.
\bibitem{dreamprover} Y.~Zhang et al., \emph{DreamProver: Evolving
Transferable Lemma Libraries via a Wake--Sleep Theorem-Proving Agent}, 2026,
\href{https://doi.org/10.48550/arXiv.2604.26311}{arXiv}.
\bibitem{hep} I.~Takahara and T.~Mizoguchi, \emph{Toward Auditable AI
Scientists: A Hypothesis Evolution Protocol for LLM Agents}, 2026,
\href{https://doi.org/10.48550/arXiv.2607.09195}{arXiv}.
\bibitem{symposium} D.~Pratt, \emph{Symposium: Trust via Auditable Records for
Communities of AI Scientist Agents}, 2026,
\href{https://doi.org/10.48550/arXiv.2608.19511}{arXiv}.
\bibitem{eurekagent} A.~Xin et al., \emph{EurekAgent: Agent Environment
Engineering is All You Need for Autonomous Scientific Discovery}, 2026,
\href{https://doi.org/10.48550/arXiv.2606.13662}{arXiv}.
\end{thebibliography}
\end{multicols}
}
\end{document}
